\documentclass[17pt]{beamer}

\usepackage[utf8]{inputenc}
\usepackage[turkish]{babel}
\usepackage[T1]{fontenc}

\usepackage{amsmath}
\usepackage{amssymb}
\usepackage{amsthm}
\usepackage{enumerate}
\usepackage{tikz}
\usepackage{transparent}
\usepackage{xcolor}

\usetheme{Warsaw}
\usecolortheme{default}
\definecolor{nigdeyesili_acik}{RGB}{3, 150, 166}
\definecolor{nigdeyesili_koyu}{RGB}{143, 209, 217}
\setbeamercolor{structure}{fg=nigdeyesili_acik}

\author[Dr. Mehmet CANEVİ]{Arş.~Gör.~Dr.~M.~Canevi\inst{1}}

\institute{
    \inst{1}%
    Bilgisayar Mühendisliği\\
    Mühendislik Fakültesi
}
    
\date[2025] {Ders Notları, Ocak 2025}
    
\logo{\transparent{0.4}\includegraphics[width = 20mm]{logo}}
    
\title[Ders 6] {Struct Yapılar ve Makrolar}
\begin{document}
%%%%%%%%%%%%%%%%%%%%%%%%%%%%%%%%%%%%%%%%%%%%%%%%%%%%%%%%%%%%%%%%%%%%%%%%%%%%%%%%
\frame{\titlepage}
\begin{frame}[fragile]{İçidekiler}
    \tableofcontents
\end{frame}
%%%%%%%%%%%%%%%%%%%%%%%%%%%%%%%%%%%%%%%%%%%%%%%%%%%%%%%%%%%%%%%%%%%%%%%%%%%%%%%%
\section{Struct yapısı}
\begin{frame}[fragile]{Struct(Yapı) nedir?}
    \begin{lstlisting}
        void fonksiyon(char* isim,char* soyisim,int yas,int cinsiyet);
    \end{lstlisting}
    fonksiyonu
    \begin{lstlisting}
        struct kisi{
            char* isim;
            char* soyisim;
            int yas;
            int cinsiyet;
        };
        void fonksiyon(struct kisi k);\end{lstlisting}
    olarak da yazılabilir.
\end{frame}
%%%%%%%%%%%%%%%%%%%%%%%%%%%%%%%%%%%%%%%%%%%%%%%%%%%%%%%%%%%%%%%%%%%%%%%%%%%%%%%%
\begin{frame}[fragile]{Struct(Yapı) nedir?}
    Böylece
    \begin{lstlisting}
        struct kisi{
            char* isim;
            char* soyisim;
            int yas;
            int cinsiyet;
        };
        struct kisi fonksiyon();\end{lstlisting}
    yazılabilir.
\end{frame}
%%%%%%%%%%%%%%%%%%%%%%%%%%%%%%%%%%%%%%%%%%%%%%%%%%%%%%%%%%%%%%%%%%%%%%%%%%%%%%%%
\begin{frame}[fragile]{Struct(Yapı) nedir?}
    Basit yapı taşlarından oluşturulan karmaşık veri tipine \lstinline{struct} denir.
    \begin{lstlisting}
        struct yapi_adi{
            yapi_tasi1;
            yapi_tasi2;
            ...
        };
        struct yapi_adi degisken_adi;\end{lstlisting}
\end{frame}
%%%%%%%%%%%%%%%%%%%%%%%%%%%%%%%%%%%%%%%%%%%%%%%%%%%%%%%%%%%%%%%%%%%%%%%%%%%%%%%%
\section{Makrolar}
\begin{frame}[fragile]{define}
    Kodları kontrol etmek için kullanılır. Programın içeriğine etki etmezler.
    \begin{lstlisting}
        #define DEGISKEN\end{lstlisting}
\end{frame}
%%%%%%%%%%%%%%%%%%%%%%%%%%%%%%%%%%%%%%%%%%%%%%%%%%%%%%%%%%%%%%%%%%%%%%%%%%%%%%%%
\begin{frame}[fragile]{ifndef}
    \begin{lstlisting}
        #ifndef 1
            float a;
        #else
            double a;
        #endif\end{lstlisting}
    \lstinline{if-else} yapısının makro karşılığıdır.
\end{frame}
%%%%%%%%%%%%%%%%%%%%%%%%%%%%%%%%%%%%%%%%%%%%%%%%%%%%%%%%%%%%%%%%%%%%%%%%%%%%%%%%
\begin{frame}[fragile]{Makro fonksiyon}
    \begin{lstlisting}
        #define yazdir(str) printf(str)\end{lstlisting}
    ile fonksiyon tanımı yapılabilmektedir.
\end{frame}
%%%%%%%%%%%%%%%%%%%%%%%%%%%%%%%%%%%%%%%%%%%%%%%%%%%%%%%%%%%%%%%%%%%%%%%%%%%%%%%%
\section{Uygulama}
\begin{frame}[fragile]{Soru}
    \begin{alertblock}{Soru-1}
        Vertex struct'ı iki boyutta koordinat verisi içermektedir. Üçgen yapısı oluşturunuz. Üçgenin çevresini hesaplayan fonksiyonu yazınız.
    \end{alertblock}
    \begin{alertblock}{Soru-2}
        Matris yapısı yazınız.
    \end{alertblock}
\end{frame}
%%%%%%%%%%%%%%%%%%%%%%%%%%%%%%%%%%%%%%%%%%%%%%%%%%%%%%%%%%%%%%%%%%%%%%%%%%%%%%%%
\end{document}